%! Author = lili
%! Date = 8/1/26

\skript{Although sum-of-pairs (SP) multiple sequence alignment is NP-complete, it is sometimes desirable to
compute optimal solutions for a moderate number of sequences; either as a sub-procedure of other
methods, or for comparison and benchmarking of heuristic alignment algorithms.

Therefore here we will come back to the original model of sum-of-pairs multiple alignment and look at
it a bit more closely. We first formally describe the multi-dimensional generalization of the pairwise
dynamic programming algorithm. Then we show how to reduce the search space considerably according to
an idea of Carrillo and Lipman.

In this section we give the formulation for distance minimization. Naturally, equivalent algorithms
exist for score maximization.}

\section{The Exact Algorithm Revisited}\label{sec:the-exact-algorithm-revisited}

\skript{The problem to be solved is the following. Given $k$ sequences $s_1, s_2, \ldots, s_k$ of lengths
    $n_1, n_2, \ldots, n_k$, respectively, the goal is to find a multiple alignment $A$ such that
    \[
        \mathrm{cost}_{SP}(A) := \sum_{1\le p\le q\le k}\mathrm{cost}_2(\pi_{p,q}(A))
    \]
    is minimum among all multiple alignments of $s_1, s_2, \ldots, s_k$.}

\subsection{The Basic Algorithm}\label{subsec:the-basic-algorithm}
\skript{We have already seen that the SP multiple alignment problem can be solved by applying the Universal
Alignment Algorithm from Chapter~5 of the ``Sequence Analysis 1'' lecture to a $k$-dimensional
alignment matrix, one dimension for each of the input sequences:

For each cell $v$ of the alignment matrix, compute in an appropriate order the dissimilarity value
    \begin{equation}
        D(v) = \min_{u\in\mathrm{pred}(v)}\bigl\{ D(u) + \mathrm{cost}(u\to v)\bigr\}. \tag{8.1}\label{eq:equation}
    \end{equation}
    More explicitly, this can be written as follows:
    \[
        D(0, 0, \ldots, 0) = 0
    \]
    \[
        D(\underbrace{i_1, i_2, \ldots, i_k}_{v}) =
        \min_{\substack{\Delta_1,\ldots,\Delta_k\in\{0,1\}\\ \Delta_1+\ldots+\Delta_k\ne 0}}
        \left\{
            D(\underbrace{i_1-\Delta_1, i_2-\Delta_2, \ldots, i_k-\Delta_k}_{\text{predecessor } u})
        + D_{SP}\underbrace{\begin{pmatrix}\Delta_1 s_1[i_1]\\ \Delta_2 s_2[i_2]\\ \vdots\\
        \Delta_k s_k[i_k]\end{pmatrix}}_{\text{alignm.\ col.}}
        \right\}
    \]
    where for a character $c\in\Sigma$ the notation $\Delta c$ denotes $\Delta c = c$ if $\Delta = 1$ and
    $\Delta c = -$ if $\Delta = 0$.}

\paragraph{Space and time complexity.} \skript{The space complexity $O(n^k)$ is clearly given by the size of the multidimensional alignment matrix,
    where $n\ge n_1, n_2, \ldots, n_k$.

    The time complexity $O(n^k\cdot 2^k\cdot k^2)$ is also easy to see in this notation: For each of the
    $O(n^k)$ cells, all $2^k-1$ predecessors have to be evaluated, given by the $2^k$ possible settings of
    the vector $(\Delta_1, \ldots, \Delta_k)$ of binary variables $\Delta_i$, without the all-zero vector
    where $\Delta_1 + \ldots + \Delta_k = 0$. In addition, the sum-of-pairs cost of the resulting
    alignment column has to be computed, which requires additional $O(k^2)$ time per predecessor.}

\subsection{Variations of the Basic Algorithm}\label{subsec:variations-of-the-basic-algorithm}

\paragraph{Memory reduction.}  \skript{From the pairwise alignment, we remember that it is possible to reduce the space for optimal alignment
    computation from quadratic to linear using Hirschberg's divide-and-conquer approach (Section~\ref{sec:hirschbergs-algorithm}).

    In principle, the same idea can also be followed for multiple alignment, but the effect is much less
    dramatic. The space reduction is by one dimension, from $O(n^k)$ to $O(n^{k-1})$. While for $k = 3$
    this has quite some effect, the space savings become less and less impressive as $k$ grows. For
    example, for $k = 12$ sequences the space consumption using this technique is still $O(n^{11})$.
    % TODO
    Figure~8.1 sketches on the left side the first division step for the pairwise case, and on the right
    side the corresponding step in the multiple (here 3-dimensional) case.}

% TODO Figure

\paragraph{Free end gaps.} \skript{In the same spirit as in semi-global pairwise alignment, it is sometimes desired to penalize gaps at
the beginning and at the end of a multiple alignment less than internal gaps (or not at all, known as
\emph{free end-gap} alignment), in order to avoid that shorter sequences are ``spread'' over longer
ones, just to get a slightly higher score that is not counter-penalized by gap costs, because these
have to be imposed anyway to account for the difference in sequence length.

This can be done in multiple alignment as in the pairwise case. It is just a different initialization
of the base cases, i.e.\ the outer edges of the $k$-dimensional alignment matrix.}

\paragraph{Affine gap costs.} \skript{Affine gap costs can be defined in a straightforward manner as a simple generalization of the pairwise
case, called \emph{natural gap costs}: The cost of a multiple alignment is the sum of all costs of the
pairwise projections, where the cost of a pairwise projection is computed with affine gap costs. In
principle it is possible to perform the computation of such alignments similarly as it is done for
pairwise alignments. However, the number of ``history matrices'' ($V$ and $H$ in the pairwise case)
    grows exponentially with the number of sequences, so that very much space is needed.

    That is why Altschul (1989) suggested to use an approximation of the exact case, so-called
    \emph{quasi-natural gap costs}. The idea is to look back by only one position in the edit graph and
    decide from the previous column in the alignment if a new gap is started in the column to be computed
    or not. The choices, compared to the ``correct'' choices in natural gap costs, are given in the
    following table whose first line depicts the different possible scenarios of pairwise alignment
    projections where an asterisk (\texttt{*}) refers to a character, a dash (\texttt{-}) refers to a
    blank, and a question mark (\texttt{?}) refers to either a character or a blank. The entries correspond
    to the cost ($d$ is gap open cost, $e$ is gap extension cost), where ``no'' means no new gap is opened
    (substitution cost):
    \begin{center}
        \setlength{\tabcolsep}{6pt}
        \begin{tabular}{l|cccccc}
            & $\substack{\texttt{?*}\\\texttt{?*}}$ & $\substack{\texttt{?-}\\\texttt{?-}}$
            & $\substack{\texttt{**}\\\texttt{*-}}$ & $\substack{\texttt{**}\\\texttt{--}}$
            & $\substack{\texttt{-*}\\\texttt{*-}}$ & $\substack{\texttt{-*}\\\texttt{--}}$\\
            \hline
            natural       & no & 0 & $d$ & $e$ & $d$ & $d/e$\\
            quasi-natural & no & 0 & $d$ & $e$ & $d$ & $d$\\
        \end{tabular}
    \end{center}
    The last case in natural gap costs cannot be decided if only one previous column is given. It would be
    $d$ for $\substack{\texttt{*---*}\\\texttt{*----}}$ and $e$ for
    $\substack{\texttt{***--}\\\texttt{**----}}$.

    With this heuristic, the overall computation becomes slower by a factor of $2^k$, yielding a time
    complexity for the complete multiple alignment computation of $O(n^k 2^{2k} k^2)$.

    The space requirements do not change in the worst case, they remain $O(n^k)$.}

\section{Carrillo and Lipman’s Search Space Reduction}\label{sec:carrillo-and-lipmans-search-space-reduction}

\skript{Carrillo and Lipman (1988) suggested an effective way to prune the search space for multiple
alignments. Their algorithm is still exponential in the worst case, but in many cases the time of
alignment computation is considerably reduced compared to the full dynamic programming algorithm from
Section~8.1, making it applicable in practice for data sets of moderate size (7--10 protein sequences
of length 300--400).

Let sequences $s_1, s_2, \ldots, s_k$ be given. The algorithm of Carrillo and Lipman is a
\emph{branch-and-bound} running time heuristic, based on a simple observation.}

\begin{idee}
    \skript{We have seen that the pairwise projections of an optimal multiple alignment are not necessarily the
    optimal pairwise alignments (making the problem hard). However, intuitively, they can also not be
    particularly bad alignments: If all pairwise projections were bad, the whole multiple alignment would
    also have a bad sum-of-pairs score.

    If we can quantify this idea, we can restrict the search space to those vertices that are part of
    close-to-optimal pairwise alignments for each pairwise projection.}
\end{idee}

\paragraph{Bounding pairwise alignment costs.} \skript{We assume that we already know some (suboptimal) multiple alignment $A^{\mathrm{heur}}$ of the given
sequences with cost $\mathrm{cost}_{SP}(A^{\mathrm{heur}})\ge\mathrm{cost}_{SP}(A^{\mathrm{opt}})$, for
example computed by a progressive alignment method or by the center-star heuristic (Section~9.3).

Remember the definition of the sum-of-pairs multiple alignment cost:
    \[
        \mathrm{cost}_{SP}(A) = \sum_{p<q}\mathrm{cost}_2(\pi_{\{p,q\}}(A)).
    \]
    We now pick a particular index pair $(x,y)$, $x<y$, and remove it from the sum. In the remaining
    pairs, we use the fact that the pairwise projections of the optimal multiple alignment cannot have a
    lower cost than the optimal corresponding pairwise alignments:
    $\mathrm{cost}_2(\pi_{\{p,q\}}(A^{\mathrm{opt}}))\ge d(s_p, s_q)$. Thus
    \[
        \begin{aligned}
            \mathrm{cost}_{SP}(A^{\mathrm{heur}}) &\ge \mathrm{cost}_{SP}(A^{\mathrm{opt}})\\
            &= \sum_{p<q}\mathrm{cost}_2(\pi_{\{p,q\}}(A^{\mathrm{opt}}))\\
            &= \mathrm{cost}_2(\pi_{\{x,y\}}(A^{\mathrm{opt}})) +
            \sum_{\substack{p<q\\(p,q)\ne(x,y)}}\mathrm{cost}_2(\pi_{\{p,q\}}(A^{\mathrm{opt}}))\\
            &\ge \mathrm{cost}_2(\pi_{\{x,y\}}(A^{\mathrm{opt}})) +
            \sum_{\substack{p<q\\(p,q)\ne(x,y)}} d(s_p, s_q)
        \end{aligned}
    \]
    for any pair $(x,y)$, $1\le x<y\le k$.

    This implies the following upper bound for the cost of the projection of an optimal multiple alignment
    on rows $x$ and $y$:
    \[
        \mathrm{cost}_2(\pi_{\{x,y\}}(A^{\mathrm{opt}})) \le
        \mathrm{cost}_{SP}(A^{\mathrm{heur}}) - \sum_{\substack{p<q\\(p,q)\ne(x,y)}} d(s_p, s_q)
        =: U_{(x,y)}.
    \]}

\begin{notte}
    \skript{Note that in order to compute the upper bound $U_{(x,y)}$, only pairwise optimal alignments and a
    heuristic multiple alignment need to be calculated. This can be done efficiently.}
\end{notte}

\paragraph{Pruning the alignment matrix.} \skript{We now keep only those cells $(i_1, \ldots, i_k)$ of the alignment matrix that satisfy the following
condition for all pairs $(x,y)$, $1\le x<y\le k$:
    \begin{center}
        \textit{The best pairwise alignment of $s_x$ and $s_y$ through cell $(i_x, i_y)$ has cost $\le U_{(x,y)}$.}
    \end{center}
    How do we know the cost of the best pairwise alignment through any given cell $(i_x, i_y)$? This
    question was answered in Section~6.1: We use the forward-backward technique. In other words, for every
    pair $(x,y)$, $x<y$, we compute the total cost matrix $T_{(x,y)}$ such that $T_{(x,y)}(i_x, i_y)$
    contains the cost of the best alignment of $s_x$ and $s_y$ that passes through the cell $(i_x, i_y)$.

    We can define the set of cells in the multi-dimensional alignment matrix that satisfy the constraint
    imposed by the pair $(x,y)$ as follows:
    \[
        B_{(x,y)} := \{(i_1, \ldots, i_k) : T_{(x,y)}(i_x, i_y)\le U_{(x,y)}\}
    \]
    From the above considerations it is clear that the projection of an SP-optimal multiple alignment on
    sequences $s_x$ and $s_y$ can not have a cost higher than $U_{(x,y)}$, implying that regions $(i,j)$
    with values $T_{(x,y)}(i,j) > U_{(x,y)}$ cannot contain such a projection.

    Thus an optimal multiple alignment passes only through cells that are in $B_{(x,y)}$ for all pairs
    $(x,y)$, i.e., in
    \[
        B := \bigcap_{x<y} B_{(x,y)}.
    \]
    This means: Instead of computing the distances in all $O(n^k)$ cells in the alignment matrix, we need
    to compute the distances of only $|B|$ cells. The whole procedure is conceptually illustrated in
    Figure~8.2. % TODo FIG

    The size of $B$ depends on many things, e.g.\ on
    \begin{itemize}
        \item the quality of the heuristic alignment with cost $\mathrm{cost}_{SP}(A^{\mathrm{heur}})$ that
        serves as an upper bound of the optimal alignment cost. The closer
        $\mathrm{cost}_{SP}(A^{\mathrm{heur}})$ is to the true optimal value, the smaller is $B$;
        \item the difference between the optimal pairwise alignment costs and the costs of the pairwise
        projections of the optimal multiple alignment.
    \end{itemize}}

% TODO FIG

\skript{Generally, if all $k$ sequences are similar and the gap costs are small, then a good heuristic
alignment (maybe even an optimal one) is easy to find, so $B$ may contain almost no additional
vertices besides the true optimal multiple alignment path. Thus, for a medium number of similar
sequences, the simple Carrillo-Lipman heuristic may already perform very well (because the relevant
regions become very small) while for even fewer but less closely related sequences even the best known
heuristics will not finish within weeks of computation. Sequence similarity plays a much higher role
than sequence length or the number of sequences.

Some implementations of the Carrillo-Lipman bound attempt to reduce the size of $B$ further by
``cheating'': Remember that we need a heuristic alignment to obtain an upper bound
    $\mathrm{cost}_{SP}(A^{\mathrm{heur}}) =: \delta$ on the optimal cost. If we have reason to suspect
    that we have found a bad alignment, and believe that a better one exists, we may replace $\delta$ by a
    smaller value, which could even be chosen differently for each pair $(x,y)$: $\delta - \epsilon_{(x,y)}$.
    Of course, since we cannot be sure that such a better alignment exists, we may reduce the size of $B$
    too much and in fact exclude the optimal multiple alignment from $B$. In practice, however, this kind
    of cheating works quite well if done carefully, and further reduces the running time.}

\paragraph{Algorithm.} \skript{We give a few remarks about implementing the above idea. We emphasize that the set of relevant matrix
cells $B$ is conceptual and does not need to be constructed explicitly. Instead, it is sufficient to
precompute all $T_{(x,y)}$ matrices, the optimal pairwise alignment scores $d(s_p, s_q)$ for all
    $(p,q)$, $1\le p<q\le k$, and a heuristic alignment $A^{\mathrm{heur}}$ to obtain $\delta$.

    It is advantageous not to implement the ``pull-type'' recurrence (8.1), but instead to use a
    ``push-type'' algorithm as explained in the following.

    Imagine that every cell has a value of $\infty$ at the beginning of the algorithm, although this
    information is not stored explicitly. Instead, a list (e.g.\ a queue or priority queue) of ``active''
    cells is maintained, that initially contains only $(0, 0, \ldots, 0)$.

    The algorithm then consists of repeatedly ``visiting'' a cell from the queue until the final cell
    $(n_1, n_2, \ldots, n_k)$ is reached. When a cell is visited, it already contains the correct alignment
    distance.

    Visiting a cell $v$ consists of the following steps.
    \begin{enumerate}
        \item We look at the $2^k-1$ successors $w$ of $v$ in the alignment matrix and check for each of them
        if it belongs to the set $B$. This is the case if $w$ is already in the queue, or if the current
        distance value plus edge cost $v\to w$ satisfies the Carrillo-Lipman bound for every index pair
        $(x,y)$. In the former case, the distance (and backpointer) information of $w$ is updated if the new
        distance value coming from $v$ is lower than the current one. In the latter case $w$ is ``activated''
        with the appropriate distance and backpointer information and added to the queue.
        \item We remove $v$ from the queue. If we want to create an alignment in the end, we need to store
        $v$'s information somewhere else to reconstruct the traceback information. If we only want the distance
        value, $v$ can now be completely discarded.
        \item We pick the next cell $x$ to visit. The algorithm works correctly if $x$ is
        \begin{itemize}
            \item a vertex of minimum distance; then the algorithm is essentially Dijkstra's algorithm for
            single-source shortest paths in the (reduced) alignment graph, or
            \item any vertex that has no active predecessors, so we can be sure that we do not need to update
            $x$'s distance information after visiting $x$.
        \end{itemize}
        Since often more than one vertex is available for visiting, a good target-driven choice may speed up
        the algorithm.
    \end{enumerate}}

\paragraph{Implementations.}
\skript{Implementations of the Carrillo-Lipman heuristic are the programs MSA (Gupta et al., 1995) and QALIGN
    (Sammeth et al., 2003b).}


